\section{Point Machine Branch}
\label{sec:methodology:point-machine-branch}

The PointMachineBranch module is responsible for validating point machine operations and ensuring that any requested movement between NORMAL and REVERSE positions is carried out safely. Point machines serve as critical mechanical switches that guide trains across diverging tracks, and their reliability is essential to prevent derailments or route conflicts. Before a position change is permitted, the module evaluates a range of safety checks, confirming that the machine exists and is active, that the requested position differs from the current one, and that the device is not mechanically or time-locked. It further verifies that all guarding track segments are unoccupied, associated protective signals display red, and no conflicts exist with neighboring points or active routes. Only if all conditions are satisfied will the operation be approved.

The module receives input such as the point machine identifier, its requested state, and real-time operational data retrieved through the DatabaseManager. It also references metadata linking the point to track topology and checks signal states to enforce protective locking. Integrated directly into the InterlockingRuleEngine, the PointMachineBranch is invoked whenever a route request involves a point movement, whether in live operation or simulation scenarios.

Diagnostics are embedded throughout the validation process. Each stage is logged with precise timestamps, reasons for approval or rejection, and references to affected entities. These logs support both real-time troubleshooting and forensic analysis of interlocking behavior, making it easier to trace the root causes of violations or timing conflicts. For instance, if point machine PM07 is requested to move from NORMAL to REVERSE, the module ensures it is active, not already in REVERSE, not locked, and that the connected track segments are clear before returning an ALLOWED result. If any condition fails, a BLOCKED response is generated along with a detailed explanation, reinforcing both operational safety and system transparency.